\section{The three-arrow case}\label{sec:three-arrow}
In this section all colored gradings use $m=3$.

\subsection{The limit at slope one}\label{sec:central}

\begin{corollary}\label{cor:central}
Let $((a_n,b_n))_{n\ge1}$ be a sequence in $\cA_3$ satisfying
\[
 \lim_{n\to\infty}a_n=\infty,\qquad
 \lim_{n\to\infty}\frac{b_n}{a_n}=1.
\]
Then,
\[
 \lim_{n\to\infty}\frac{\log E_3(a_n,b_n)}{a_n}
 =\log\left(\frac{256}{27}\right).
\]
\end{corollary}
\begin{proof}
Let us reduce the dimension pair to $(d_n,d_n+k_n)$.
For $n\ge1$, define
\[
 d_n=\min(a_n,b_n),\qquad k_n=|b_n-a_n|.
\]
Since $(a_n,b_n)\in\cA_3$, we have $\gcd(a_n,b_n)=1$.
If $a_n=b_n$, then $a_n=b_n=1$.
The condition $a_n\to\infty$ therefore gives $n_0\in\Z_{>0}$ such that $k_n\ge1$ for $n\ge n_0$.
For $n\ge n_0$, we have
\[
 \gcd(d_n,k_n)=\gcd(\min(a_n,b_n),|b_n-a_n|)=\gcd(a_n,b_n)=1.
\]
Claim~\ref{c:moduli-symmetries-identities} of Lemma~\ref{lem:moduli-symmetries} gives
\[
 E_3(a_n,b_n)=E_{k_n}(d_n).
\]

Let us compute the limits of the new dimensions.
The hypothesis $\frac{b_n}{a_n}\to1$ gives
\[
 \frac{d_n}{a_n}=\min\left(1,\frac{b_n}{a_n}\right)\longrightarrow1,
 \qquad
 \frac{k_n}{d_n}
 =\frac{\left|\frac{b_n}{a_n}-1\right|}{\min\left(1,\frac{b_n}{a_n}\right)}
 \longrightarrow0.
\]
In particular, $d_n\to\infty$.
Define the two index sets
\[
 J_1=\{n\ge n_0:k_n=1\},\qquad
 J_2=\{n\ge n_0:k_n\ge2\}.
\]

First, let us compute the limit along $J_1$ when $J_1$ is infinite.
Equation~\eqref{eq:adjacent-three} and Stirling's formula give
\[
 \log E_1(d_n)=d_n\log\beta-\frac52\log d_n+O(1)
 \qquad(n\to\infty,\ n\in J_1).
\]
Dividing by $d_n$ gives
\[
 \lim_{\substack{n\to\infty\\n\in J_1}}
 \frac{\log E_{k_n}(d_n)}{d_n}=\log\beta.
\]

Second, let us compute the limit along $J_2$ when $J_2$ is infinite.
Enumerate $J_2=\{n_1<n_2<\cdots\}$.
The sequence $((d_{n_j},k_{n_j}))_{j\ge1}$ satisfies
\[
 \begin{gathered}
 k_{n_j}\ge2,\qquad \gcd(d_{n_j},k_{n_j})=1\qquad(j\ge1),\\
 \lim_{j\to\infty}d_{n_j}=\infty,\qquad
 \lim_{j\to\infty}\frac{k_{n_j}}{d_{n_j}}=0.
 \end{gathered}
\]
Claim~\ref{c:central-lower-limit} of Theorem~\ref{thm:central-lower} and
Corollary~\ref{cor:central-delete-limsup} therefore give
\[
 \lim_{j\to\infty}\frac{\log E_{k_{n_j}}(d_{n_j})}{d_{n_j}}=\log\beta.
\]

Let us combine the two cases.
We have $J_1\sqcup J_2=\{n\ge n_0\}$.
For each infinite $J_i$, $i\in[2]_{\Z}$, the limit of $\frac{\log E_{k_n}(d_n)}{d_n}$ along $n\in J_i$ is $\log\beta$. Each finite $J_i$ contributes only finitely many terms.
Thus,
\[
 \frac{\log E_3(a_n,b_n)}{a_n}
 =\frac{d_n}{a_n}\frac{\log E_{k_n}(d_n)}{d_n}
 \longrightarrow1\cdot\log\beta=\log\left(\frac{256}{27}\right).
\]
\end{proof}

\subsection{Adjacent exchanges}
\begin{definition}[Integer lists and balance operations]\label{def:balance-operations}
Let $\operatorname{IntList}=\bigsqcup_{n\in\Z_{\ge0}}\Z^n$.
\begin{enumerate}[font=\itshape]
\item For $L=(a_1,\ldots,a_n)\in\operatorname{IntList}$, let $|L|=n$.
For $M=(b_1,\ldots,b_m)\in\operatorname{IntList}$, define
\[
 L\ccn M=(a_1,\ldots,a_n,b_1,\ldots,b_m).
\]
For $a\in\Z$ and $s\in\Z_{\ge0}$, define
\[
 \operatorname{Repeat}(a,s)=\underbrace{(a,\ldots,a)}_{s\text{ entries}},\qquad
 \operatorname{Repeat}(a,0)=().
\]
\item We define the set of \textit{balance states} by
$\operatorname{BalState}=\Z\times\operatorname{IntList}$.
\item\label{d:horizontal-balance-step} For $s\in\Z_{\ge0}$, define
\[
 \operatorname{HStep}_s:\operatorname{BalState}\longrightarrow\operatorname{BalState},\qquad
 \operatorname{HStep}_s(b,L)=(b+1,\operatorname{Repeat}(b+1,s)\ccn L).
\]
\item\label{d:vertical-balance-step} For $h\in\Z_{\ge0}$, define
\[
 \mathcal V_h=\{(b,(a)\ccn L_0)\in\operatorname{BalState}:a\le b-h\}
\]
and
\[
 \operatorname{VStep}_h:\mathcal V_h\longrightarrow\operatorname{BalState},\qquad
 \operatorname{VStep}_h(b,(a)\ccn L_0)=(b-h,L_0).
\]
We call $\operatorname{HStep}_s$ and $\operatorname{VStep}_h$ the horizontal and vertical balance operations.
\item\label{d:balance-list-order} For $\delta\in\Z_{\ge0}$ and integer lists $L$ and $L'$, define
\begin{equation}\label{eq:exchange-list-comparison}
 L'\preceq_\delta L
 \ \Longleftrightarrow\
 |L'|=|L|\ \text{ and }\ L'_i\le L_i+\delta\quad(i\in[|L|]_{\Z}).
\end{equation}
\item\label{d:boundary-balance-list} Let $Y$ be a tNE path. Let $g\in\operatorname{CtdG}(Y)$. Let $r\in[0,|Y|]_{\Z}$.
List the horizontal edges of $Y[0,r]$ as $e_1,\ldots,e_h$ in path order.
For $i\in[h]_{\Z}$, let
\[
 b_i=\operatorname{eBal}_{Y,g}(e_i),\qquad
 d_i=\operatorname{Debt}_{Y,g}(e_i,r).
\]
Define the \textit{boundary balance list} by
\[
 \operatorname{BalList}_{Y,g}(r)
 =\operatorname{Repeat}(b_h,d_h)\ccn\cdots\ccn\operatorname{Repeat}(b_1,d_1).
\]
For $h=0$, this list is $()$.
\end{enumerate}
\end{definition}

\begin{lemma}[Comparison of balance operations]\label{lem:balance-list-comparison}
Let $\nu\in\Z_{\ge0}$ and $\delta\in\Z_{\ge0}$.
For $t\in[\nu]_{\Z}$, let $\varepsilon_t\in\{E,N\}$ and $q_t\in\Z_{\ge0}$.
Assume the following.
\begin{hypotheses}
\item\label{hyp:balance-initial} $z_0=(b_0,L_0)\in\operatorname{BalState}$ and $L'_0\preceq_\delta L_0$.
Let $z'_0=(b_0+\delta,L'_0)$.
\item\label{hyp:balance-sequence} The states $z_t=(b_t,L_t)$ for $t\in[\nu]_{\Z}$ satisfy
\[
 z_t=\begin{cases}
 \operatorname{HStep}_{q_t}(z_{t-1}),&\varepsilon_t=E,\\
 \operatorname{VStep}_{q_t}(z_{t-1}),&\varepsilon_t=N.
 \end{cases}
\]
For every $t\in[\nu]_{\Z}$ with $\varepsilon_t=N$, we have $z_{t-1}\in\mathcal V_{q_t}$.
\end{hypotheses}
Starting from $z'_0$, the states $z'_t=(b'_t,L'_t)$ are well-defined by
\[
 z'_t=\begin{cases}
 \operatorname{HStep}_{q_t}(z'_{t-1}),&\varepsilon_t=E,\\
 \operatorname{VStep}_{q_t}(z'_{t-1}),&\varepsilon_t=N
 \end{cases}
 \qquad(t\in[\nu]_{\Z}).
\]
For every $t\in[0,\nu]_{\Z}$,
\[
 b'_t=b_t+\delta,\qquad L'_t\preceq_\delta L_t.
\]
In particular, $z'_{t-1}\in\mathcal V_{q_t}$ whenever $\varepsilon_t=N$.
\end{lemma}
\begin{proof}
Let us prove the assertions by induction on $t$.
The case $t=0$ is assumption~\ref{hyp:balance-initial}.
Suppose $b'_{t-1}=b_{t-1}+\delta$ and $L'_{t-1}\preceq_\delta L_{t-1}$.
If $\varepsilon_t=E$, the horizontal operation increases each balance by one and prepends $q_t$ entries.
The new entries are $b_{t-1}+1$ and $b_{t-1}+1+\delta$, respectively.
Thus, $b'_t=b_t+\delta$ and $L'_t\preceq_\delta L_t$.
If $\varepsilon_t=N$, assumption~\ref{hyp:balance-sequence} gives
$L_{t-1}=(a)\ccn M$ and $a\le b_{t-1}-q_t$.
The equal list lengths give $L'_{t-1}=(a')\ccn M'$.
By the induction hypothesis,
\[
 a'\le a+\delta\le b_{t-1}-q_t+\delta=b'_{t-1}-q_t.
\]
Hence, $z'_{t-1}\in\mathcal V_{q_t}$.
Both vertical operations delete the first entry and subtract $q_t$ from the balance.
Thus, $b'_t=b_t+\delta$ and $L'_t\preceq_\delta L_t$.
\end{proof}

\Needspace{19\baselineskip}
\begin{definition}[Adjacent exchange in a shell]\label{def:shell-adjacent-exchange}
Define $\operatorname{AdjShellExch}$ to be the set of tuples
\[
 \xi=(P,Q,o,u,v,A,B)
\]
such that $P$ and $Q$ are NE paths, $o\in\Z^2$, $u\in[2]_{\Z}$, $v\in[2]_{\Z}$, and
$(A,B)\in\{(E\ccn N,N\ccn E),(N\ccn E,E\ccn N)\}$.
We call an element of this set an \textit{adjacent exchange in a shell}.
Let $\xi=(P,Q,o,u,v,A,B)\in\operatorname{AdjShellExch}$.
\begin{enumerate}[font=\itshape]
\item\label{d:adj-shell-exch-paths} Define
\[
\begin{aligned}
 X_\xi&=\operatorname{Trans}(E\ccn P\ccn A\ccn Q,o),& X'_\xi&=\operatorname{Trans}(E\ccn P\ccn B\ccn Q,o),\\
 Y_\xi&=\operatorname{Shell}(X_\xi;u,v,0),&Y'_\xi&=\operatorname{Shell}(X'_\xi;u,v,0).
\end{aligned}
\]
\item\label{d:adj-shell-exch-edges} Let $Y=Y_\xi$ and $Y'=Y'_\xi$. Let $(e_i)_{i\in[|Y|_E]_{\Z}}$ and $(n_j)_{j\in[|Y|_N]_{\Z}}$
be the horizontal and vertical edge lists of $Y$, respectively, in path order.
Let $(e'_i)_{i\in[|Y'|_E]_{\Z}}$ and $(n'_j)_{j\in[|Y'|_N]_{\Z}}$
be the horizontal and vertical edge lists of $Y'$, respectively, in path order.
For each $j\in[|Y'_\xi|_N]_{\Z}$, let $p_j\in[|Y'_\xi|]_{\Z}$ be the unique index satisfying $(Y'_\xi)_{p_j}=n'_j$.
\item\label{d:adj-shell-edge-map} Define the bijection
\[
 \tau_\xi:\operatorname{NE}(Y_\xi)\longrightarrow\operatorname{NE}(Y'_\xi),\qquad
 \tau_\xi(e_i)=e'_i,\quad \tau_\xi(n_j)=n'_j.
\]
\item\label{d:adj-shell-exchanged-edges} Define
\[
 r_\xi=u+2+|P|,\qquad
 i_\xi=|P|_E+3,\qquad j_\xi=u+|P|_N+1,\qquad
 x_\xi=e_{i_\xi},\quad w_\xi=n_{j_\xi}.
\]
\item\label{d:adj-shell-gradings} Define
\[
 \mathcal G_\xi=
 \{f\in\operatorname{CtdG}_3(X_\xi):\operatorname{Surp}(X_\xi,f)\ge0\}.
\]
For $f\in\mathcal G_\xi$, define
\[
 g_{\xi,f}=\operatorname{GShell}_f(X_\xi;u,v,0),\qquad
 \overline g_{\xi,f}=g_{\xi,f}\circ\tau_\xi^{-1}.
\]
\item\label{d:adj-shell-counter} Let $f\in\mathcal G_\xi$ and let $s=g_{\xi,f}(x_\xi)$.
For $s>0$, define integers $q_i$ for $i\in[r_\xi+2,|Y'_\xi|]_{\Z}$ by
\[
 q_{r_\xi+2}=s-1,\qquad
 q_i=\begin{cases}
 q_{i-1}+\overline g_{\xi,f}((Y'_\xi)_i),&(Y'_\xi)_i\in E(Y'_\xi),\\
 \max(q_{i-1}-1,0),&(Y'_\xi)_i\in N(Y'_\xi)
 \end{cases}
 \quad(i>r_\xi+2).
\]
Define
\[
 \mathcal J_{\xi,f}=\{j\in[|Y'_\xi|_N]_{\Z}:p_j>r_\xi+2,\ q_{p_j-1}=0\}.
\]
For $s=0$, let $\mathcal J_{\xi,f}=\varnothing$.
\item\label{d:adj-shell-repair-index} Define the \textit{vertical repair index} by
\[
 J_\xi(f)=\begin{cases}
 \min\mathcal J_{\xi,f},&A=E\ccn N,\ s>0,\ \mathcal J_{\xi,f}\ne\varnothing,\\
 j_\xi,&\text{otherwise}.
 \end{cases}
\]
\item\label{d:adj-shell-repair-grading} Let $J=J_\xi(f)$ and define
\[
\begin{gathered}
 h=g_{\xi,f}(w_\xi),\qquad h_J=\overline g_{\xi,f}(n'_J),\\
 \vartheta=\operatorname{bBal}_{Y_\xi,g_{\xi,f}}(r_\xi)+1,\qquad
 b_J=\operatorname{bBal}_{Y'_\xi,\overline g_{\xi,f}}(p_J-1).
\end{gathered}
\]
Define the \textit{grade reduction} by
\[
 D_\xi(f)=\begin{cases}
 h-\max(h-1,0),&A=E\ccn N,\ s=0,\\
 0,&A=N\ccn E,\ s=0,\\
 h,&A=N\ccn E,\ s>0,\\
 h_J-\min\{h_J,\max(b_J-\vartheta,0)\},&A=E\ccn N,\ s>0.
 \end{cases}
\]
Define the \textit{repaired grading} $\operatorname{RepGr}_\xi(f)$ on $Y'_\xi$ by
\[
 \operatorname{RepGr}_\xi(f)(e'_i)=g_{\xi,f}(e_i),\qquad
 \operatorname{RepGr}_\xi(f)(n'_j)=
 \begin{cases}
 g_{\xi,f}(n_j)-D_\xi(f),&j=J_\xi(f),\\
 g_{\xi,f}(n_j),&j\ne J_\xi(f).
 \end{cases}
\]
\end{enumerate}
\end{definition}

\begin{lemma}[The vertical repair index]\label{lem:exchange-repair-index}
Let $\xi=(P,Q,o,u,v,A,B)\in\operatorname{AdjShellExch}$.
\begin{enumerate}
\item\label{c:repair-index-exists} If $f\in\mathcal G_\xi$, $A=E\ccn N$, and $g_{\xi,f}(x_\xi)>0$, then $\mathcal J_{\xi,f}\ne\varnothing$.
The edge $n'_{J_\xi(f)}$ lies in the copy of $X'_\xi$ in $Y'_\xi$.
\item\label{c:repair-index-horizontal} If $f\in\mathcal G_\xi$, $\widetilde f\in\mathcal G_\xi$, and
\[
 g_{\xi,f}(e_i)=g_{\xi,\widetilde f}(e_i)
 \qquad(i\in[|Y_\xi|_E]_{\Z}),
\]
then $J_\xi(f)=J_\xi(\widetilde f)$.
\end{enumerate}
\end{lemma}
\begin{proof}
\prf{Claim~\ref{c:repair-index-exists}} Let us prove $\mathcal J_{\xi,f}\ne\varnothing$.
Let $g=g_{\xi,f}$, $s=g(x_\xi)>0$, $r=r_\xi$, and $T=u+1+|X_\xi|$.
The shell grading $g=g_{\xi,f}$ belongs to $\operatorname{CtdG}_3(Y_\xi)$ by Claim~\ref{c:shell-connected} of Lemma~\ref{lem:shell}.
Since $f$ is connected, the last vertical edge of $X_\xi$ has h-precursor $(X_\xi)_1$.
Claim~\ref{c:shell-copy-precursor} of Lemma~\ref{lem:shell} and Claim~\ref{c:assigned-prefix} of Lemma~\ref{lem:precursor-prefix-counts} give
\[
 \operatorname{Debt}_{Y_\xi,g}((Y_\xi)_{u+2},r)\ge1.
\]
Since $g$ is connected, Claim~\ref{c:assigned-prefix} of Lemma~\ref{lem:precursor-prefix-counts} gives $\operatorname{Debt}_{Y_\xi,g}((Y_\xi)_1,r)\ge1$.
Let
\[
 \operatorname{BalList}_{Y_\xi,g}(r)=(a)\ccn L_0.
\]
The entry contributed by $(Y_\xi)_1$ follows an entry contributed by $(Y_\xi)_{u+2}$.
Hence, $L_0$ contains the entry contributed by $(Y_\xi)_1$, and $|L_0|\ge1$.
The edge $w_\xi$ immediately follows $x_\xi$, and $g(x_\xi)=s>0$.
Definition~\ref{def:stack}(1) gives $P_{Y_\xi,g}(w_\xi)=x_\xi$.
The node balance condition~\eqref{eq:connected-edge-balance} gives
\[
 \operatorname{eBal}_{Y_\xi,g}(x_\xi)
 \le\operatorname{eBal}_{Y_\xi,g}(w_\xi)
 =\operatorname{eBal}_{Y_\xi,g}(x_\xi)-g(w_\xi).
\]
Hence, $g(w_\xi)=0$.
Definition~\ref{def:shell-adjacent-exchange}(\ref{d:adj-shell-counter}) gives $q_t\ge0$ for $t\in[r+2,|Y'_\xi|]_{\Z}$.
By Definition~\ref{def:balance-operations}(6),
\[
 |\operatorname{BalList}_{Y_\xi,g}(r+2)|
 =s+|L_0|=q_{r+2}+1+|L_0|.
\]

Let us prove that $q_{p_j-1}=0$ for some $p_j\in[r+3,T]_{\Z}$.
Suppose, for a contradiction, that $q_{p_j-1}>0$ for every vertical position $p_j\in[r+3,T]_{\Z}$.
For $t\in[r+3,T]_{\Z}$, we have $\tau_\xi((Y_\xi)_t)=(Y'_\xi)_t$.
At a horizontal edge, both $q_t$ and $|\operatorname{BalList}_{Y_\xi,g}(t)|$ increase by $g((Y_\xi)_t)$.
At a vertical edge at position $t=p_j$, the inequality $q_{t-1}>0$ gives $q_t=q_{t-1}-1$, and Definition~\ref{def:stack}(1) decreases $|\operatorname{BalList}_{Y_\xi,g}(t)|$ by $1$.
Induction on $t$ gives
\[
 |\operatorname{BalList}_{Y_\xi,g}(t)|=q_t+1+|L_0|
 \qquad(t\in[r+2,T]_{\Z}).
\]

Let us apply the shell boundary identity at $T$.
Claim~\ref{c:shell-boundary-debts} of Lemma~\ref{lem:shell}, the equality $\operatorname{eBal}_{Y_\xi,g}((Y_\xi)_1)=1$, and Definition~\ref{def:balance-operations}(6) give $\operatorname{BalList}_{Y_\xi,g}(T)=(1)$.
Therefore,
\[
 1=|\operatorname{BalList}_{Y_\xi,g}(T)|
 =q_T+1+|L_0|\ge2,
\]
a contradiction.
Thus, some $j\in\mathcal J_{\xi,f}$ satisfies $p_j\le T$.
Definition~\ref{def:shell-adjacent-exchange}(\ref{d:adj-shell-repair-index}) gives $r+2<p_{J_\xi(f)}\le T$.
Hence, $n'_{J_\xi(f)}$ lies in the copy of $X'_\xi$ in $Y'_\xi$.

\prf{Claim~\ref{c:repair-index-horizontal}} Let us check which grades determine the repair index.
The common horizontal grades determine $s$ and every horizontal increment in the recursion for $q_i$.
The vertical recursion $q_i=\max(q_{i-1}-1,0)$ does not use a vertical grade.
Hence, $\mathcal J_{\xi,f}=\mathcal J_{\xi,\widetilde f}$ and $J_\xi(f)=J_\xi(\widetilde f)$.
\end{proof}

\begin{lemma}[The repaired grading]\label{lem:uncolored-adjacent-exchange}
Let $\xi\in\operatorname{AdjShellExch}$ and $f\in\mathcal G_\xi$.
Let
\[
 g=g_{\xi,f},\qquad g'=\operatorname{RepGr}_\xi(f),\qquad
 j=J_\xi(f),\qquad\delta=D_\xi(f).
\]
Then,
\begin{enumerate}
\item\label{c:uncolored-exchange-grades} $g'\in\operatorname{CtdG}_3(Y'_\xi)$ and $0\le\delta\le g(n_j)\le3$.
\item\label{c:uncolored-exchange-surplus} We have
\[
 \operatorname{Surp}(Y'_\xi,g')=\operatorname{Surp}(X_\xi,f)+\delta.
\]
\end{enumerate}
\end{lemma}
\begin{proof}
\prf{Claim~\ref{c:uncolored-exchange-grades}} Let us compare the h-precursors before and after exchanging the two edges.
Let $\xi=(P,Q,o,u,v,A,B)$ and $(X,X',Y,Y')=(X_\xi,X'_\xi,Y_\xi,Y'_\xi)$.
Put $c=\operatorname{Surp}(X,f)$.
Claim~\ref{c:shell-connected} of Lemma~\ref{lem:shell} gives $g\in\operatorname{CtdG}_3(Y)$.
The formulas for $D_\xi(f)$ give $0\le\delta\le g(n_j)\le3$.
Thus, $g'\in G_3(Y')$.
For $r\in[0,|Y|]_{\Z}$, let
\[
 b=\operatorname{bBal}_{Y,g}(r),\qquad L=\operatorname{BalList}_{Y,g}(r).
\]
For $r<|Y|$, Definition~\ref{def:stack}(1) and the node balance condition in
Definition~\ref{def:connected-grading}(1) give the operations $\operatorname{HStep}_{g(Y_{r+1})}$ when $Y_{r+1}\in E(Y)$ and $\operatorname{VStep}_{g(Y_{r+1})}$ when $Y_{r+1}\in N(Y)$.
Claim~\ref{c:node-boundary-balances} of Lemma~\ref{lem:node-structure}
implies that every entry of $L$ is at most $b$.
Claims~\ref{c:shell-boundary-debts} and~\ref{c:shell-boundary-balances} of Lemma~\ref{lem:shell} and Definition~\ref{def:balance-operations}(6) give
\[
\begin{aligned}
 (\operatorname{bBal}_{Y,g}(u+1),\operatorname{BalList}_{Y,g}(u+1))&=(1,(1)),\\
 (\operatorname{bBal}_{Y,g}(u+1+|X|),\operatorname{BalList}_{Y,g}(u+1+|X|))&=(c+2,(1)).
\end{aligned}
\]
The entry $1$ in $L=(1)$ is contributed by $Y_1$.
The vertical edges of $X$ remove every entry that $X$ contributes.

Let us construct states $z'_t=(b'_t,L'_t)$ for $t\in[0,|Y'|]_{\Z}$.
Let $z'_0=(0,())$ and use the recursion
\[
 z'_t=\begin{cases}
 \operatorname{HStep}_{g'(Y'_t)}(z'_{t-1}),&Y'_t\in E(Y'),\\
 \operatorname{VStep}_{g'(Y'_t)}(z'_{t-1}),&Y'_t\in N(Y').
 \end{cases}
\]

Let us verify $z'_{t-1}\in\mathcal V_{g'(Y'_t)}$ before the exchanged pair.
For $t\in[0,r_\xi]_{\Z}$, we have
\[
 z'_t=(\operatorname{bBal}_{Y,g}(t),\operatorname{BalList}_{Y,g}(t)).
\]

Let us compare the balance states at boundaries $r_\xi$ and $r_\xi+2$.
Let $\tau=\tau_\xi$, $x=x_\xi$, and $w=w_\xi$.
Let $s=g(x)$, $h=g(w)$, and $w'=\tau(w)$.
At boundary $r_\xi$ of $Y$, let $(b,(a)\ccn L_0)$ be the balance state.
The list $L_0$ is nonempty: the first edge of $Y$ contributes its entry $1$. The first edge of the copy of $X$ has positive debt before the exchanged pair.
By the definition of $g'=\operatorname{RepGr}_\xi(f)$,
\[
 g'(\tau(e_i))=g(e_i)\quad(i\in[|Y|_E]_{\Z}),\qquad
 g'(\tau(n_i))=g(n_i)\quad(i\in[|Y|_N]_{\Z}\setminus\{J_\xi(f)\}).
\]
At boundary $r_\xi+2$, let
\[
\begin{aligned}
 (b_{\mathrm{old}},L_{\mathrm{old}})&=
 (\operatorname{bBal}_{Y,g}(r_\xi+2),\operatorname{BalList}_{Y,g}(r_\xi+2)),\\
 (b_{\mathrm{new}},L_{\mathrm{new}})&=z'_{r_\xi+2}.
\end{aligned}
\]

First, let us verify the balance condition when $A=E\ccn N$ and $s=0$.
We have $J_\xi(f)=j_\xi$, $g'(w')=\max(h-1,0)$, and $\delta=h-\max(h-1,0)$.
Validity of $\operatorname{VStep}_h$ at $w$ gives $a\le b+1-h$.
We also have $a\le b$.
Hence,
\[
 a\le b-\max(h-1,0).
\]
Thus,
\[
 L_{\mathrm{new}}=L_{\mathrm{old}},\qquad
 b_{\mathrm{new}}=b_{\mathrm{old}}+\delta,\qquad \delta\in[0,1]_{\Z}.
\]

Second, let us verify the balance condition when $A=N\ccn E$ and $s=0$.
We have $\delta=0$ and $g'(w')=h$.
The inequality $a\le b-h$ implies $a\le b+1-h$.
Thus, $(b_{\mathrm{new}},L_{\mathrm{new}})=(b_{\mathrm{old}},L_{\mathrm{old}})$.

Third, let us verify the balance condition when $A=N\ccn E$ and $s>0$.
We have $J_\xi(f)=j_\xi$, $g'(w')=0$, and $\delta=h$.
The states at boundary $r_\xi+2$ satisfy
\[
 b_{\mathrm{old}}=b-h+1,\qquad
 L_{\mathrm{old}}=\operatorname{Repeat}(b-h+1,s)\ccn L_0,
\]
\[
 b_{\mathrm{new}}=b+1,\qquad
 L_{\mathrm{new}}=\operatorname{Repeat}(b+1,s-1)\ccn(a)\ccn L_0.
\]
Since $a\le b-h$, we have $L_{\mathrm{new}}\preceq_h L_{\mathrm{old}}$ by~\eqref{eq:exchange-list-comparison}.

Fourth, let us verify the balance condition when $A=E\ccn N$ and $s>0$.
Validity of $\operatorname{VStep}_h$ at $w$ immediately after $x$ gives $b+1\le b+1-h$. Thus, $h=0$.
Let $\vartheta=b+1$.
The states at boundary $r_\xi+2$ satisfy
\[
 L_{\mathrm{old}}=\operatorname{Repeat}(\vartheta,s-1)\ccn(a)\ccn L_0,
\]
\[
 L_{\mathrm{new}}=\operatorname{Repeat}(\vartheta,s)\ccn L_0,\qquad
 b_{\mathrm{new}}=b_{\mathrm{old}}=\vartheta.
\]

Let us locate the entry whose grade requires a change in the fourth case.
Mark entry $s$ of $L_{\mathrm{new}}$.
Claim~\ref{c:repair-index-exists} of Lemma~\ref{lem:exchange-repair-index} gives
\[
 r_\xi+2<p_j\le u+1+|X|.
\]
The edge $T=\tau^{-1}(n'_j)$ removes the marked entry.
At boundary $t\in[r_\xi+2,p_j-1]_{\Z}$, the mark has position $q_t+1$, where $q_t$ is the counter in Definition~\ref{def:shell-adjacent-exchange}(\ref{d:adj-shell-counter}).

Let us prove that the operations before $T$ are valid.
For $t\in[r_\xi+2,p_j-1]_{\Z}$, we prove by induction that there is a list $H_t$ satisfying
\[
\begin{aligned}
 \operatorname{BalList}_{Y,g}(t)&=H_t\ccn(a)\ccn L_0,\\
 L'_t&=H_t\ccn(\vartheta)\ccn L_0,\\
 |H_t|&=q_t,\qquad (H_t)_i\ge\vartheta\quad(i\in[q_t]_{\Z}),\\
 b'_t&=\operatorname{bBal}_{Y,g}(t)\ge\vartheta.
\end{aligned}
\]
At $t=r_\xi+2$, take $H_t=\operatorname{Repeat}(\vartheta,s-1)$.

First, let us prove the induction step when $Y'_{t+1}\in E(Y')$ and $t+1<p_j$.
Both balances increase by one.
Every entry inserted by this horizontal edge equals $b'_t+1\ge\vartheta$.
Set $H_{t+1}=\operatorname{Repeat}(b'_t+1,g'(Y'_{t+1}))\ccn H_t$.
The identity $q_{t+1}=q_t+g'(Y'_{t+1})$ verifies $|H_{t+1}|=q_{t+1}$.
Both fixed tails, $(a)\ccn L_0$ and $(\vartheta)\ccn L_0$, are unchanged.

Second, let us prove the induction step when $Y'_{t+1}\in N(Y')$ and $t+1<p_j$.
The definition of $j=J_\xi(f)$ gives $q_t>0$.
Let $a_t$ be the common first list entry.
The operation $\operatorname{VStep}_{g(Y_{t+1})}$ on $Y$ at position $t+1$ is valid, so
\[
 b'_t-g'(Y'_{t+1})
 =\operatorname{bBal}_{Y,g}(t)-g(Y_{t+1})
 \ge a_t\ge\vartheta.
\]
The inequality $a_t\le b'_t-g'(Y'_{t+1})$ verifies the condition in Definition~\ref{def:balance-operations}(\ref{d:vertical-balance-step}) for the operation at $Y'_{t+1}$.
Since $q_t>0$, deleting the first entry of $H_t$ gives $H_{t+1}$ of length $q_{t+1}=q_t-1$.
Both fixed tails are unchanged, and $b'_{t+1}=b'_t-g'(Y'_{t+1})\ge\vartheta$.

Let us verify the vertical operation at $\tau(T)$.
Let $b_T$ be the boundary balance immediately before $T$ and let $q=g(T)$.
We have $b_T\ge\vartheta$.
At $T'=\tau(T)=n'_{J_\xi(f)}$, the definition of $D_\xi(f)$ gives
\[
 g'(T')=k:=\min\{q,b_T-\vartheta\},\qquad \delta=q-k.
\]
Then, $0\le k\le q$ and $b_T-k\ge\vartheta$.
Since $q_{p_j-1}=0$, the lists immediately before $T$ and $T'$ are $(a)\ccn L_0$ and $(\vartheta)\ccn L_0$, respectively.
Deleting their first entries gives
\[
 L'_{p_j}=\operatorname{BalList}_{Y,g}(p_j)=L_0,\qquad
 b'_{p_j}=\operatorname{bBal}_{Y,g}(p_j)+\delta.
\]

Let us verify connectedness.
Lemma~\ref{lem:balance-list-comparison} gives boundary balance $c+2+\delta$ and a list of length one at $u+1+|X|$.
The entry contributed by $Y'_1$ equals $1$ and stays last in every nonempty list.
The original lists are nonempty throughout the copy of $X$.
Equality of list lengths therefore prevents this entry from being removed.
Hence,
\[
 z'_{u+1+|X|}=(c+2+\delta,(1)).
\]
The grades on the final $R_v$ are $v-1$ on its horizontal edge, zero on its first
$v-1$ vertical edges, and two on its last edge.
After its horizontal edge the state is
\[
 (c+3+\delta,\operatorname{Repeat}(c+3+\delta,v-1)\ccn(1)).
\]
The zero-grade vertical operations are valid.
The last vertical operation is valid because $1\le c+1+\delta$.
Therefore,
\[
 z'_{|Y'|}=(c+1+\delta,()).
\]

Let us define the h-precursor from the lists $L'_t$.
Assign each list entry to the horizontal edge that inserted it.
For $t\in[|Y'|]_{\Z}$ with $Y'_t\in N(Y')$, let $Y'_{r_t}$ be the horizontal edge assigned to the first entry of $L'_{t-1}$.
The list is nonempty because $z'_{t-1}\in\mathcal V_{g'(Y'_t)}$.
The order of the list gives
\[
 r_t=\max\left\{i\in[t-1]_{\Z}:
 \begin{array}{l}
 Y'_i\in E(Y'),\\
 g'(Y'_i)>\#\{j\in[t-1]_{\Z}:Y'_j\in N(Y'),\ r_j=i\}
 \end{array}\right\}.
\]
Thus, $P_{Y',g'}(Y'_t)=Y'_{r_t}$ satisfies Definition~\ref{def:stack}(1).

Let us verify the node balance and the node at $Y'_1$.
The first entry of $L'_{t-1}$ is $\operatorname{eBal}_{Y',g'}(Y'_{r_t})$.
Membership in $\mathcal V_{g'(Y'_t)}$ gives
\[
 \operatorname{eBal}_{Y',g'}(Y'_{r_t})
 \le b'_{t-1}-g'(Y'_t)
 =\operatorname{eBal}_{Y',g'}(Y'_t)
 \qquad(Y'_t\in N(Y')).
\]
Every list entry is removed because $L'_{|Y'|}=()$.
The entry from $Y'_1$ is removed at $Y'_{|Y'|}$.
Thus,
\[
 \operatorname{Node}_{Y',g'}(Y'_1)=Y',\qquad g'\in\operatorname{CtdG}_3(Y').
\]

\prf{Claim~\ref{c:uncolored-exchange-surplus}} Let us compute the surplus change.
Only the vertical grade at $n'_j$ changes by $-\delta$. Hence,
\[
 \operatorname{Surp}(Y',g')=c+\delta\in[c,c+3]_{\Z}.
\]
\end{proof}

\begin{corollary}[The adjacent-exchange repair map]\label{cor:adjacent-repair-map}
Let $\xi\in\operatorname{AdjShellExch}$.
Then, the formula
\[
 \Theta_\xi(f)=(\operatorname{RepGr}_\xi(f),J_\xi(f),D_\xi(f))
\]
defines a map
\[
 \Theta_\xi:\mathcal G_\xi\longrightarrow
 \operatorname{CtdG}_3(Y'_\xi)\times[|Y_\xi|_N]_{\Z}\times[0,3]_{\Z}.
\]
\end{corollary}
\begin{proof}
Let us verify the target of $\Theta_\xi$.
Fix $f\in\mathcal G_\xi$.
Claim~\ref{c:uncolored-exchange-grades} of Lemma~\ref{lem:uncolored-adjacent-exchange} gives
$\operatorname{RepGr}_\xi(f)\in\operatorname{CtdG}_3(Y'_\xi)$ and $D_\xi(f)\in[0,3]_{\Z}$.
The definition of $J_\xi(f)$ gives $J_\xi(f)\in[|Y_\xi|_N]_{\Z}$.
\end{proof}


\Needspace{25\baselineskip}
\begin{lemma}[An adjacent edge exchange inside a shell]\label{lem:exchange-precursors}
Let $\xi=(P,Q,o,u,v,A,B)\in\operatorname{AdjShellExch}$.
Let
\[
 (X,X',Y,Y')=(X_\xi,X'_\xi,Y_\xi,Y'_\xi).
\]
Let $e_1,\ldots,e_{|Y|_E}$ and $n_1,\ldots,n_{|Y|_N}$ be the horizontal and vertical edges of $Y$ in path order.
Let $e'_1,\ldots,e'_{|Y'|_E}$ and $n'_1,\ldots,n'_{|Y'|_N}$ be the horizontal and vertical edges of $Y'$ in path order.
There are maps
\[
\begin{aligned}
 \Psi&:\bigcup_{c\ge0}\mathcal C_c(X)\longrightarrow
             \bigcup_{c\ge0}\mathcal C_c(Y'),\\
 \kappa&:\bigcup_{c\ge0}\mathcal C_c(X)\longrightarrow
 \{(h,d):h\in[0,3]_{\Z},\ d\in[\tbinom3h]_{\Z}\}
\end{aligned}
\]
with the following properties.
For each $(f,\gamma)\in\bigcup_{c\ge0}\mathcal C_c(X)$, let $g=\operatorname{GShell}_{f}(X;u,v,0)$.
Let $\eta$ be the shell coloring in~\eqref{eq:shell-recolor} and let
$(g',\gamma')=\Psi(f,\gamma)$.
Then, we have the following.

\begin{enumerate}
\item\label{c:adjacent-exchange-surplus} We have
\[
 \operatorname{Surp}(X,f)\le\operatorname{Surp}(Y',g')
 \le\operatorname{Surp}(X,f)+3.
\]
\item\label{c:adjacent-exchange-horizontal} For $i\in[|Y|_E]_{\Z}$, we have
\[
 (g'(e'_i),\gamma'(e'_i))=(g(e_i),\eta(e_i)).
\]
\item\label{c:adjacent-exchange-vertical} We have
\[
 \#\{j\in[|Y|_N]_{\Z}:(g'(n'_j),\gamma'(n'_j))\ne(g(n_j),\eta(n_j))\}\le1.
\]
\item\label{c:adjacent-exchange-recovery} The map
\[
 (f,\gamma)\longmapsto\bigl(\Psi(f,\gamma),\kappa(f,\gamma)\bigr)
\]
is injective.
\end{enumerate}
\end{lemma}
\begin{proof}
\prf{Claim~\ref{c:adjacent-exchange-surplus}} Let us define $\Psi$ and $\kappa$.
Fix $(f,\gamma)\in\mathcal C_c(X)$.
For every $e\in\operatorname{NE}(X)$, the condition $\gamma(e)\in[\binom3{f(e)}]_{\Z}$ gives $f(e)\le3$. Thus, $f\in\mathcal G_\xi$.
Use the map $\Theta_\xi$ from Corollary~\ref{cor:adjacent-repair-map} and let $(g'_f,j_f,\delta_f)=\Theta_\xi(f)$.
Let $\tau=\tau_\xi$. Define
\[
 \gamma'=\operatorname{Recol}(Y,Y',\tau,g'_f;g,\eta),\qquad
 \Psi(f,\gamma)=(g'_f,\gamma'),\qquad
 \kappa(f,\gamma)=(g(n_{j_f}),\eta(n_{j_f})).
\]
Claim~\ref{c:uncolored-exchange-surplus} of Lemma~\ref{lem:uncolored-adjacent-exchange} gives the surplus bounds.

\prf{Claim~\ref{c:adjacent-exchange-horizontal}} Let us check the horizontal grade and color pairs.
Definition~\ref{def:shell-adjacent-exchange}(\ref{d:adj-shell-repair-grading}) keeps every horizontal grade.
Equation~\eqref{eq:color-update} therefore keeps every horizontal color.

\prf{Claim~\ref{c:adjacent-exchange-vertical}} Let us check how many vertical grade and color pairs change.
The repair in Definition~\ref{def:shell-adjacent-exchange}(\ref{d:adj-shell-repair-grading}) changes at most the vertical grade at $n'_{j_f}$.
Equation~\eqref{eq:color-update} keeps the colors at all other vertical edges.

\prf{Claim~\ref{c:adjacent-exchange-recovery}} Let us recover the input from the output and the one changed pair.
The pair $(g',\gamma')$ determines $g(e_i)=g'(e'_i)$ for every $i\in[|Y|_E]_{\Z}$.
Claim~\ref{c:repair-index-horizontal} of Lemma~\ref{lem:exchange-repair-index}
therefore determines $j_f$ from those grades.
For $j\in[|Y|_N]_{\Z}\setminus\{j_f\}$, Claim~\ref{c:adjacent-exchange-vertical} and the definition of $\gamma'$ give $(g'(n'_j),\gamma'(n'_j))=(g(n_j),\eta(n_j))$.
Recover $(g(n_{j_f}),\eta(n_{j_f}))=\kappa(f,\gamma)$.
Reverse the exchange and remove the two outer simple runs.
The result is $(f,\gamma)$. Thus, $(\Psi(f,\gamma),\kappa(f,\gamma))$ determines $(f,\gamma)$.
\end{proof}

\begin{definition}[Exchange decompositions]\label{def:exchange-decomposition}
Let $X$ be a tNE path. If $X=\emptyset$, let $\operatorname{ExcDecomp}(X)=\emptyset$.
Suppose $X\ne\emptyset$. An \textit{exchange decomposition} of $X$ is a tuple
$(B^{\mathrm L},S_1,S_2,B^{\mathrm R})$ of NE paths such that $S_1$ and $S_2$ are simple runs and
\[
 X=\operatorname{Trans}(B^{\mathrm L}\ccn S_1\ccn S_2\ccn B^{\mathrm R},\operatorname{Start}_{\mathrm{pt}}(X)),\qquad
 |S_1|_N\ne|S_2|_N.
\]
We denote the set of these tuples by $\operatorname{ExcDecomp}(X)$.
\end{definition}

\begin{lemma}\label{lem:exchange}
Let $X$ be a binary run. Let $u_0\in[2]_{\Z}$ and $v_0\in[2]_{\Z}$.
Let $(B^{\mathrm L},S_1,S_2,B^{\mathrm R})\in\operatorname{ExcDecomp}(X)$.
Let $X'=\operatorname{Exc}(X;|B^{\mathrm L}|,S_1,S_2)$.
Let $Y=R_{u_0}\ccn X\ccn R_{v_0}$ and $Y'=R_{u_0}\ccn X'\ccn R_{v_0}$.
There is a map
\[
 \Phi:\bigcup_{c\ge0}\mathcal C_c(X)\longrightarrow
       \bigcup_{c\ge0}\mathcal C_c(Y')
\]
with the following properties.
For each $(f,\gamma)\in\bigcup_{c\ge0}\mathcal C_c(X)$, let $g=\operatorname{GShell}_{f}(X;u_0,v_0,0)$
and define $\eta$ by~\eqref{eq:shell-recolor}.
Let $(g',\gamma')=\Phi(f,\gamma)$.
Label the horizontal and vertical edges of $Y$ by $e_i$ and $n_j$, respectively.
Label the horizontal and vertical edges of $Y'$ by $e'_i$ and $n'_j$, respectively, in path order.
Then, we have the following.

\begin{enumerate}
\item\label{c:exchange-connected} The grading $g'$ satisfies the following.
\begin{enumerate}[label=(\alph*),ref=\theenumi(\alph*)]
\item\label{c:exchange-connectedness} $g'\in\operatorname{CtdG}(Y')$.
\item\label{c:exchange-residual} We have
\[
 \operatorname{Surp}(X,f)\le\operatorname{Surp}(Y',g')
 \le\operatorname{Surp}(X,f)+3.
\]
\end{enumerate}
\item\label{c:exchange-horizontal} We have
\[
 (g'(e'_i),\gamma'(e'_i))=(g(e_i),\eta(e_i))
 \quad(i\in[|Y|_E]_{\Z}).
\]
\item\label{c:exchange-vertical} We have
\[
 \#\{j\in[|Y|_N]_{\Z}:(g'(n'_j),\gamma'(n'_j))\ne(g(n_j),\eta(n_j))\}\le1.
\]
\item\label{c:exchange-fibers} For every $(h,\delta)\in\bigcup_{c\ge0}\mathcal C_c(Y')$, we have
\[
 \#\Phi^{-1}(\{(h,\delta)\})\le8.
\]
\end{enumerate}
\end{lemma}
\begin{proof}
\prf{Claim~\ref{c:exchange-connectedness}} Let us specify the tuple in Definition~\ref{def:shell-adjacent-exchange}.
Define
\[
\begin{gathered}
 o=\operatorname{Start}_{\mathrm{pt}}(X),\qquad H=B^{\mathrm L}\ccn E\ccn N,\\
 P=\operatorname{Trans}(H[1,|H|],(-1,0)),\qquad Q=N\ccn B^{\mathrm R},\\
 (A,B)=
 \begin{cases}
 (E\ccn N,N\ccn E),&(S_1,S_2)=(R_1,R_2),\\
 (N\ccn E,E\ccn N),&(S_1,S_2)=(R_2,R_1).
 \end{cases}
\end{gathered}
\]
Since $X$ is binary, $H$ begins with $E$.
Thus, $P$ and $Q$ are NE paths and $E\ccn P=H$.
It follows that
\[
 \xi=(P,Q,o,u_0,v_0,A,B)\in\operatorname{AdjShellExch},\qquad
 (X_\xi,X'_\xi,Y_\xi,Y'_\xi)=(X,X',Y,Y').
\]
Apply Lemma~\ref{lem:exchange-precursors} to $\xi$ and let $\Phi=\Psi$.
The inclusion $\Psi(f,\gamma)\in\bigcup_{c\ge0}\mathcal C_c(Y')$ gives $g'\in\operatorname{CtdG}(Y')$.

\prf{Claim~\ref{c:exchange-residual}}
Claim~\ref{c:adjacent-exchange-surplus} of Lemma~\ref{lem:exchange-precursors} gives the surplus bounds.

\prf{Claim~\ref{c:exchange-horizontal}}
Claim~\ref{c:adjacent-exchange-horizontal} of Lemma~\ref{lem:exchange-precursors} preserves every horizontal grade--color pair.

\prf{Claim~\ref{c:exchange-vertical}}
Claim~\ref{c:adjacent-exchange-vertical} of Lemma~\ref{lem:exchange-precursors} bounds the number of changed vertical grade--color pairs by one.

\prf{Claim~\ref{c:exchange-fibers}} Let us bound each fiber by the number of recovery choices.
By Claim~\ref{c:adjacent-exchange-recovery} of Lemma~\ref{lem:exchange-precursors},
the map $\kappa$ is injective on each fiber of $\Phi$.
The set of possible values of $\kappa$ has cardinality
\[
 \sum_{h=0}^{3}\binom3h=8.
\]
Hence, every fiber of $\Phi$ has cardinality at most eight.
\end{proof}
\begin{definition}[Wrapped exchange]\label{def:wrapped-exchange}
Let $X$ be a binary run. Let $u_0\in[2]_{\Z}$ and $v_0\in[2]_{\Z}$.
Let $(B^{\mathrm L},S_1,S_2,B^{\mathrm R})\in\operatorname{ExcDecomp}(X)$.
Let $X'=\operatorname{Exc}(X;|B^{\mathrm L}|,S_1,S_2)$.
Let $\operatorname{WrapExch}(X,X';u_0,v_0)$ be the set of \textit{wrapped exchanges}

\[
 \Phi\colon\bigcup_{c\in\Z_{\ge0}}\mathcal C_c(X)
 \longrightarrow\bigcup_{c\in\Z_{\ge0}}\mathcal C_c(R_{u_0}\ccn X'\ccn R_{v_0})
\]
satisfying Claims~\ref{c:exchange-connected}--\ref{c:exchange-fibers}
of Lemma~\ref{lem:exchange}.
\end{definition}




\Needspace{10\baselineskip}
\subsection{Mechanical comparison by exchanges}
\begin{definition}\label{def:binary-run-count}
Let $W$ be a binary run with $n$ simple runs. Let
\[
 W=\operatorname{Trans}(R_{w_1}\ccn \cdots\ccn R_{w_n},\operatorname{Start}_{\mathrm{pt}}(W)).
\]
For $r\in[2]_{\Z}$, define
\[
 \operatorname{RunCount}_r(W)=\#\{i\in[n]_{\Z}:w_i=r\}.
\]
\end{definition}

\begin{definition}[Binary prefix discrepancies]\label{def:prefix-discrepancy}
Let $Z$ and $Y$ be binary runs with $n$ and $m$ simple runs, respectively. Let
\[
 Z=\operatorname{Trans}(R_{z_1}\ccn \cdots\ccn R_{z_n},\operatorname{Start}_{\mathrm{pt}}(Z)),\qquad
 Y=\operatorname{Trans}(R_{y_1}\ccn \cdots\ccn R_{y_m},\operatorname{Start}_{\mathrm{pt}}(Y)).
\]
For $j\in[0,\min(n,m)]_{\Z}$, define
\[
 \operatorname{PrefDiff}_j(Z,Y)=
 \#\{i\in[j]_{\Z}:z_i=2\}-\#\{i\in[j]_{\Z}:y_i=2\}.
\]
For $h\in[0,\min(n,m)]_{\Z}$, define
$\operatorname{PrefSum}_h(Z,Y)=\sum_{j=1}^h|\operatorname{PrefDiff}_j(Z,Y)|$.
\end{definition}

\begin{lemma}[Exchanges between binary runs]\label{lem:binary-prefix-exchanges}
Let $n\in\Z_{>0}$.
Let $\mathbf t\in\Z^2$ and let
\[
 Z=\operatorname{Trans}(R_{z_1}\ccn \cdots\ccn R_{z_n},\mathbf t),\qquad
 Y=\operatorname{Trans}(R_{y_1}\ccn \cdots\ccn R_{y_n},\mathbf t)
\]
be binary runs.
Assume $\operatorname{RunCount}_2(Z)=\operatorname{RunCount}_2(Y)$.
For $j\in[n]_{\Z}$, let
\[
 D_j=\operatorname{PrefDiff}_j(Z,Y).
\]
Then, there are $\ell\in\Z_{\ge0}$ and a run-edit list $\Omega=(\omega_s)_{s\in[\ell]_{\Z}}$ such that
\[
 \operatorname{RunEResult}(\Omega,Z)=Y.
\]
Every $\omega_s$ exchanges two adjacent simple runs of distinct vertical lengths, and
\begin{equation}\label{eq:binary-exchange-distance}
 \ell=\operatorname{PrefSum}_n(Z,Y).
\end{equation}
\end{lemma}
\begin{proof}
Let us decrease $\operatorname{PrefSum}_n(Z,Y)$ by adjacent exchanges to prove~\eqref{eq:binary-exchange-distance}.
If $Z=Y$, the empty run-edit list suffices. Suppose $Z\ne Y$ and let
\[
 a=\min\{i\in[n]_{\Z}:z_i\ne y_i\}.
\]
The entries before $a$ agree. Since the total counts of each vertical length agree, there is $i>a$ such that $z_i=y_a$.
Let
\[
 b=\min\{i\in[a+1,n]_{\Z}:z_i=y_a\}.
\]
For $j\in[a,b-1]_{\Z}$, minimality of $b$ gives $z_j=3-y_a$. Therefore,
\[
 D_j=\begin{cases}
 -\#\{i\in[a,j]_{\Z}:y_i=2\}\le-1,&y_a=2,\\
 \#\{i\in[a,j]_{\Z}:y_i=1\}\ge1,&y_a=1.
 \end{cases}
\]

Let us move $R_{z_b}=R_{y_a}$ to position $a$.
Set $Z^{(b)}=Z$. For $k=b-1,b-2,\ldots,a$, define
\[
 Z^{(k)}=\operatorname{Exc}\left(Z^{(k+1)};
 \sum_{i=1}^{k-1}(1+z_i),R_{3-y_a},R_{y_a}\right).
\]
For $k\in[a,b-1]_{\Z}$, the simple runs of $Z^{(k+1)}$ in positions $k$ and $k+1$ are $R_{3-y_a}$ and $R_{y_a}$, respectively, up to translation.
Their vertical lengths are distinct. Let
\[
 D_j^{(k)}=\operatorname{PrefDiff}_j(Z^{(k)},Y)
 \qquad(j\in[n]_{\Z},\ k\in[a,b]_{\Z}).
\]
Only the prefix ending at position $k$ changes in the exchange. Thus,
\[
 D_j^{(k)}=\begin{cases}
 D_k^{(k+1)}+1,&j=k,\ y_a=2,\\
 D_k^{(k+1)}-1,&j=k,\ y_a=1,\\
 D_j^{(k+1)},&j\ne k.
 \end{cases}
\]
The exchanges at positions greater than $k$ do not change this prefix. Thus, $D_k^{(k+1)}=D_k$.
The bounds $D_k\le-1$ when $y_a=2$ and $D_k\ge1$ when $y_a=1$ give
\[
 |D_k^{(k)}|=|D_k^{(k+1)}|-1,\qquad
 \operatorname{PrefSum}_n(Z^{(k)},Y)
 =\operatorname{PrefSum}_n(Z^{(k+1)},Y)-1.
\]

Let us prove that the sequence of exchanges terminates.
After these $b-a$ exchanges, $Z^{(a)}$ agrees with $Y$ through position $a$.
If $Z^{(a)}\ne Y$, use $Z^{(a)}$ as the new input $Z$.
Recompute its vertical lengths $z_1,\ldots,z_n$, choose $a$ and $b$ by the defining minima, and perform the exchanges for $k=b-1,\ldots,a$.
Each application increases the number of initial simple runs that agree with $Y$. Thus, the process ends after finitely many applications.
Consequently, the resulting run-edit list ends at $Y$ and has length
\[
 \ell=\operatorname{PrefSum}_n(Z,Y)-\operatorname{PrefSum}_n(Y,Y)
 =\operatorname{PrefSum}_n(Z,Y).
\]
\end{proof}

\begin{lemma}[Completion of a binary run]\label{lem:binary-completion}
Let $L\in\Z_{>0}$ and $C\in\Z_{\ge0}$.
Let $\mathbf t\in\Z^2$ and let
\[
 X=\operatorname{Trans}(R_{x_1}\ccn \cdots\ccn R_{x_L},\mathbf t),\qquad
 Y=\operatorname{Trans}(R_{y_1}\ccn \cdots\ccn R_{y_{L+C}},\mathbf t)
\]
be binary runs.
For $r\in[2]_{\Z}$, assume
\[
 b_r=\operatorname{RunCount}_r(Y)-\operatorname{RunCount}_r(X)\ge0.
\]
Let $Z=X\ccn R_1^{b_1}\ccn R_2^{b_2}$.
Then, $b_1+b_2=C$, and there are $\ell\in\Z_{\ge0}$ and a run-edit list $\Omega=(\omega_s)_{s\in[\ell]_{\Z}}$ such that
\[
 \operatorname{RunEResult}(\Omega,Z)=Y.
\]
Every $\omega_s$ exchanges two adjacent simple runs of distinct vertical lengths, and
\begin{equation}\label{eq:binary-completion-bound}
 \ell\le\operatorname{PrefSum}_L(X,Y)+\frac{C(C-1)}2.
\end{equation}
\end{lemma}
\begin{proof}
Let us check the counts after appending the missing simple runs.
The two run counts sum to the number of simple runs. Thus, $b_1+b_2=C$.
Since $Z=X\ccn R_1^{b_1}\ccn R_2^{b_2}$, we have
$\operatorname{RunCount}_r(Z)=\operatorname{RunCount}_r(X)+b_r=\operatorname{RunCount}_r(Y)$ for $r\in[2]_{\Z}$.
Apply Lemma~\ref{lem:binary-prefix-exchanges}.
For $j\in[L+C]_{\Z}$, let
\[
 D_j=\operatorname{PrefDiff}_j(Z,Y).
\]
For $j\le L$, the first $j$ simple runs of $Z$ belong to $X$.
For $j=L+r$ with $r\in[C]_{\Z}$, equality of the total counts gives
\[
 |D_{L+r}|\le C-r,\qquad
 \sum_{r=1}^{C}|D_{L+r}|\le\sum_{r=1}^{C}(C-r)=\frac{C(C-1)}2.
\]
Adding $\sum_{j=1}^L|D_j|$ gives~\eqref{eq:binary-completion-bound}.
\end{proof}

\begin{definition}[Binary run-edit lists]\label{def:binary-run-edit-lists}
Let $U$ and $V$ be binary runs and let $K\in\Z_{\ge0}$ and $k\in\Z_{\ge0}$.
For a run-edit list $\Omega=(\omega_s)_{s\in[K+k]_{\Z}}$ from $U$ to $V$, let
\[
 U_s=\operatorname{RunEResult}((\omega_1,\ldots,\omega_s),U)
 \qquad(s\in[0,K+k]_{\Z}).
\]
We define membership in $\operatorname{BinEdit}_{K,k}(U,V)$ as follows.
We have $\Omega\in\operatorname{BinEdit}_{K,k}(U,V)$ if and only if every $U_s$ is binary for $s\in[0,K+k]_{\Z}$ and there is a partition
\[
 [K+k]_{\Z}=I_R\sqcup I_X,\qquad \#I_R=K,\qquad \#I_X=k,
\]
with the following conditions.
\begin{enumerate}[font=\itshape]
\item For $s\in I_R$, there are $j_s\in[2]_{\Z}$ and $i_s\in[0,|U_{s-1}|]_{\Z}$ such that
\[
 \omega_s(U_{s-1})=\operatorname{Ins}(U_{s-1},i_s,R_{j_s}),\qquad
 i_s=|U_{s-1}|\ \text{or}\ (U_{s-1})_{i_s+1}\in E(U_{s-1}).
\]
\item For $s\in I_X$, the run-edit $\omega_s$ exchanges adjacent simple runs $A$ and $B$ such that
\[
 (|A|_N,|B|_N)\in\{(1,2),(2,1)\}.
\]
\end{enumerate}
\end{definition}

\begin{lemma}[Comparison of mechanical factors by exchanges]\label{lem:binary-factor-comparison}
Let $r\in(1,2)$ be irrational and $0<\eta<\frac14$. Let
\[
 \tau=\frac12\min(r-1,2-r),\qquad C=\ceil{\frac4\tau}.
\]
Then, there are $L\in\Z_{>0}$ and $\delta>0$ depending only on $r$ and $\eta$ with the following properties.
\begin{enumerate}
\item\label{c:binary-factor-constants} We have
\[
 C\le\eta L,\qquad \delta<\min\left\{\tau,\frac{\eta}{4L}\right\}.
\]
\item\label{c:binary-factor-uniform} For each $\sigma_1\in(r-\delta,r+\delta)$, $\sigma_2\in(r-\delta,r+\delta)$,
$u_1\in[0,1)$, $u_2\in[0,1)$, $p\in\Z_{\ge0}$, $V=\operatorname{MechRun}(\sigma_2,u_2)$, and translated copy $X$ of
$\operatorname{RunFactor}_{1,L}(\operatorname{MechRun}(\sigma_1,u_1))$,
there are $q\in\Z_{\ge0}$, $k\in\Z_{\ge0}$, $\zeta\in\Z^2$, and a run-edit list $\Omega$ with the following properties.
\begin{enumerate}[label=(\alph*),ref=\theenumi(\alph*)]
\item\label{c:binary-factor-budget} We have $C+q\le\eta L$.
\item\label{c:binary-factor-target} We have
\[
 \Omega\in\operatorname{BinEdit}_{C+q,k}
 \bigl(X,\operatorname{Trans}(\operatorname{RunFactor}_{p+1,L+C+q}(V),\zeta)\bigr).
\]
\item\label{c:binary-factor-exchanges} We have $k\le7\eta L$.
\end{enumerate}
\end{enumerate}
\end{lemma}
\begin{proof}
\prf{Claim~\ref{c:binary-factor-constants}} Let us choose a rational approximation common to both slopes.
Choose relatively prime $p_0\in\Z_{>0}$ and $q_0\in\Z_{>0}$ satisfying~\eqref{eq:mechanical-approximation}.
Choose $L\in\Z_{>0}$ and $\delta>0$ such that
\[
 q_0+C\le\eta L,\qquad C(C+1)\le\eta L,\qquad
 \delta<\min\left\{\tau,q_0^{-2},\frac{\eta}{4L}\right\}.
\]
These constants depend only on $r$ and $\eta$. They satisfy $C\le q_0+C\le\eta L$ and $\delta<\min\{\tau,\frac\eta{4L}\}$.

\prf{Claim~\ref{c:binary-factor-budget}} Let us bound the count of each simple-run length.
For $i\in[2]_{\Z}$, we have $1+\tau<\sigma_i<2-\tau$ and $|\sigma_i-\frac{p_0}{q_0}|<2q_0^{-2}$.
Let us choose $q$ and the simple runs inserted before $X$.
Apply Claim~\ref{c:phase-selection} of Lemma~\ref{lem:phase-prefix-discrepancy}
with offset $b=p$. Let $q=\min J$ and $v_q=\{u_2+(p+q)\sigma_2\}$ as in Claim~\ref{c:phase-selection} of Lemma~\ref{lem:phase-prefix-discrepancy}.
Then, $0\le q<q_0$ and $|u_1-v_q|<\eta$.
If $q>0$, insert $\operatorname{RunFactor}_{p+1,q}(V)$ immediately before $X$.
We have $q+C\le q_0+C\le\eta L$.

\prf{Claim~\ref{c:binary-factor-target}} Let us compare the completed factor with its target.
Let $Y$ be the translated copy of $\operatorname{RunFactor}_{p+q+1,L+C}(V)$
that starts at $\operatorname{Start}_{\mathrm{pt}}(X)$.

Let us specify the simple runs appended to $X$.
A mechanical factor of $a$ simple runs and slope $\sigma$ has
$a(\sigma-1)+e$ simple runs of length two and $a(2-\sigma)-e$ of length one
for some $|e|<1$.
For $j\in[2]_{\Z}$, define
\[
 b_j=\operatorname{RunCount}_j(Y)-\operatorname{RunCount}_j(X).
\]
The inequalities $\sigma_2-1>\tau$, $2-\sigma_2>\tau$, and
$|\sigma_2-\sigma_1|<2\delta$ give
\[
 b_j>\tau C-2-2\delta L>0,\qquad b_1+b_2=C.
\]
Append $b_1$ copies of $R_1$ and $b_2$ copies of $R_2$ immediately after $X$.

Let us construct the exchanges from $X\ccn R_1^{b_1}\ccn R_2^{b_2}$ to $Y$.
Lemma~\ref{lem:binary-completion} supplies a run-edit list $\mathcal S=(S_1,\ldots,S_\nu)$ from $X\ccn R_1^{b_1}\ccn R_2^{b_2}$ to $Y$
consisting of exchanges of adjacent simple runs of distinct vertical lengths.
Let $H$ be the prefix of $q$ simple runs inserted before $X$.
When a local run-edit acts on the current factor $Z$ after $H$, replace its edge index $i$ by $|H|+i$.
Replace $\operatorname{Exc}(Z;i,R_a,R_b)$ by
\[
 \operatorname{Exc}(H\ccn Z;|H|+i,R_a,R_b).
\]

Let us identify the path obtained by the complete run-edit list.
The $q$ prefix insertions followed by the $C$ appended insertions and these exchanges give $\Omega$.
Let $k=\nu$ and
\[
 \zeta=\operatorname{Start}_{\mathrm{pt}}(\operatorname{RunEResult}(\Omega,X))
       -(p,\floor{u_2+p\sigma_2}).
\]
The inserted and exchanged simple runs have vertical lengths in $[2]_{\Z}$. We have
\[
 \operatorname{RunEResult}(\Omega,X)
 =\operatorname{Trans}(\operatorname{RunFactor}_{p+1,L+C+q}(V),\zeta).
\]
Each insertion is between simple runs or at an end.
Thus, Definition~\ref{def:binary-run-edit-lists} gives the membership assertion in Claim~\ref{c:binary-factor-target}.

\prf{Claim~\ref{c:binary-factor-exchanges}} Let us bound the number of exchanges in $\mathcal S$.
For $i\in[L]_{\Z}$, the difference of the length-two prefix counts is
\[
 D_i=\operatorname{PrefDiff}_i(X,Y)=\floor{u_1+i\sigma_1}-\floor{v_q+i\sigma_2}.
\]
Claim~\ref{c:phase-prefix-mass} of Lemma~\ref{lem:phase-prefix-discrepancy} gives
\[
 \operatorname{PrefSum}_L(X,Y)=\sum_{i=1}^L|D_i|\le\frac{11}{2}\eta L.
\]
The bound in Lemma~\ref{lem:binary-completion} gives
\[
 k=\nu\le\sum_{i=1}^L|D_i|+\frac{C(C-1)}2
 \le\frac{13}{2}\eta L\le7\eta L.
\]
\end{proof}

\begin{lemma}[Completion with one exceptional simple run]\label{lem:binary-final-completion}
Let $t\in\Z_{>0}$, $C\in\Z_{>0}$, $p\in\Z_{\ge0}$, $\sigma_1\in(1,2)$, $\sigma_2\in(1,2)$, $u_1\in[0,1)$, and $u_2\in[0,1)$.
Assume
\[
 C\min(\sigma_2-1,2-\sigma_2)>3+t|\sigma_1-\sigma_2|.
\]
Let $A_i$ be simple runs satisfying $|A_i|_N\in[2]_{\Z}$ for $i\in[t]_{\Z}$.
Let $X=A_1\ccn \cdots\ccn A_t$.
Assume
\[
 |A_i|_N=\floor{u_1+i\sigma_1}-\floor{u_1+(i-1)\sigma_1}
 \qquad(i\in[t-1]_{\Z}).
\]
Let
\[
 Y=\operatorname{Trans}\bigl(\operatorname{RunFactor}_{p+1,t+C}(\operatorname{MechRun}(\sigma_2,u_2)),
  \operatorname{Start}_{\mathrm{pt}}(X)-(p,\floor{u_2+p\sigma_2})\bigr).
\]
Then, there are $k\in\Z_{\ge0}$ and a run-edit list $\Omega$ satisfying the following.
\begin{enumerate}
\item\label{c:binary-final-target} We have $\Omega\in\operatorname{BinEdit}_{C,k}(X,Y)$.
\item\label{c:binary-final-exchanges} We have $k\le(t+C)^2$.
\end{enumerate}
\end{lemma}
\begin{proof}
\prf{Claim~\ref{c:binary-final-target}} Let us specify the simple runs appended to $X$.
For $j\in[2]_{\Z}$, put $b_j=\operatorname{RunCount}_j(Y)-\operatorname{RunCount}_j(X)$.
Replacing the final mechanical simple run by $A_t$ changes either run count by at most one.
Definition~\ref{def:mechanical-paths}(\ref{d:mechanical-sequence}) gives
\[
\begin{gathered}
 \operatorname{RunCount}_2(X)=t(\sigma_1-1)+e_1,\qquad
 \operatorname{RunCount}_2(Y)=(t+C)(\sigma_2-1)+e_2,\\
 |e_1|<2,\qquad |e_2|<1.
\end{gathered}
\]
We have
\[
 \operatorname{RunCount}_1(X)=t-\operatorname{RunCount}_2(X),\qquad
 \operatorname{RunCount}_1(Y)=t+C-\operatorname{RunCount}_2(Y).
\] Hence,
\[
 b_j>C\min(\sigma_2-1,2-\sigma_2)-3-t|\sigma_1-\sigma_2|>0,
 \qquad b_1+b_2=C.
\]
Append $R_1^{b_1}\ccn R_2^{b_2}$ to $X$.
Lemma~\ref{lem:binary-prefix-exchanges} supplies exchanges from $X\ccn R_1^{b_1}\ccn R_2^{b_2}$ to $Y$.
Let $\Omega$ consist of these $C$ insertions followed by the exchanges. Let $k$ be the number of exchanges.
Every intermediate path is binary. Hence, $\Omega\in\operatorname{BinEdit}_{C,k}(X,Y)$.

Let us bound the exchanges in the completed final factor.
\prf{Claim~\ref{c:binary-final-exchanges}} Let $Z=X\ccn R_1^{b_1}\ccn R_2^{b_2}$.
For $i\in[t+C]_{\Z}$, we have $|\operatorname{PrefDiff}_i(Z,Y)|\le t+C$.
Lemma~\ref{lem:binary-prefix-exchanges} gives
\[
 k\le\operatorname{PrefSum}_{t+C}(Z,Y)\le(t+C)^2.
\]
\end{proof}

\Needspace{23\baselineskip}
\begin{lemma}\label{lem:binary-mechanical}
Let $r\in(1,2)$ be irrational and let $0<\eta<1$.
Then, there are constants $0<\delta<\min(r-1,2-r)$ and $B>0$
depending only on $r$ and $\eta$ with the following property.
For each $(2,r,\eta,\delta,B)$-admissible tuple
$(l,T,\sigma_1,\sigma_2,u_1,u_2,(A_1,\ldots,A_l))$ and
\[
 U=A_1\ccn\cdots\ccn A_l,\qquad V=\operatorname{MechRun}(\sigma_2,u_2),
\]
there are $K\in\Z_{\ge0}$, $k\in\Z_{\ge0}$, $\zeta\in\Z^2$, and a run-edit list $\Omega$ such that we have the following.
\begin{enumerate}
\item\label{c:binary-mechanical-target} $\Omega\in\operatorname{BinEdit}_{K,k}
 \bigl(U,\operatorname{Trans}(\operatorname{RunFactor}_{1,T}(V),\zeta)\bigr)$.
\item\label{c:binary-mechanical-insertions} $K=T-l$.
\item\label{c:binary-mechanical-exchanges} $k\le7\eta l+B$.
\end{enumerate}
\end{lemma}
\begin{proof}
\prf{Claim~\ref{c:binary-mechanical-target}} Let us choose the parameters for the binary factor comparison.
Let $\eta_0=\min(\eta,\frac18)$. Then, $0<\eta_0<\frac14$ and $\eta_0\le\eta$.
Choose $\tau$, $C$, $L$, and $\delta$ from Lemma~\ref{lem:binary-factor-comparison} with parameter $\eta_0$. Let
$B=C+(L+C)^2+1$.
The inequalities $|\sigma_i-r|<\delta<\min(r-1,2-r)$ give $1<\sigma_i<2$ for $i\in[2]_{\Z}$.
Thus, every simple run of $U$ and $V$ has vertical length in $[2]_{\Z}$.
Let
\[
\begin{gathered}
 h=\floor{\frac{l-1}{L}},\qquad t=l-hL\in[L]_{\Z},\\
 X_s=A_{(s-1)L+1}\ccn \cdots\ccn A_{sL}\quad(s\in[h]_{\Z}),\qquad X_*=A_{hL+1}\ccn \cdots\ccn A_l.
\end{gathered}
\]

Let us apply Claims~\ref{c:binary-factor-budget}--\ref{c:binary-factor-exchanges} of Lemma~\ref{lem:binary-factor-comparison} to $X_1,\ldots,X_h$.
Let $p_0=0$.
For $s\in[h]_{\Z}$, use the source intercept $\{u_1+(s-1)L\sigma_1\}$ and the target index $p_{s-1}$.
Claim~\ref{c:binary-factor-target} of Lemma~\ref{lem:binary-factor-comparison} gives $q_s\in\Z_{\ge0}$ and a run-edit list
$\Omega_s=(\omega_{s,1},\ldots,\omega_{s,n_s})$ from $X_s$ to a translated copy of
$\operatorname{RunFactor}_{p_{s-1}+1,L+C+q_s}(V)$.
Define $p_s=p_{s-1}+L+C+q_s$.

Let us bound the number of run-edits in each $\Omega_s$.
For $s\in[h]_{\Z}$, let $K_s$ and $k_s$ be its numbers of simple-run insertions and adjacent exchanges, respectively.
Claims~\ref{c:binary-factor-budget} and~\ref{c:binary-factor-exchanges} of Lemma~\ref{lem:binary-factor-comparison} give
\[
 K_s\le\eta_0L,\qquad k_s\le7\eta_0L\qquad(s\in[h]_{\Z}).
\]

Let us apply Claims~\ref{c:binary-final-target} and~\ref{c:binary-final-exchanges} of Lemma~\ref{lem:binary-final-completion} to $X_*$.
Use the source intercept $\{u_1+hL\sigma_1\}$ and target index $p_h$.
Since $\tau C\ge4$, $\min(\sigma_2-1,2-\sigma_2)>\tau$, and $2\delta t<\frac{\eta_0}2<1$, we have
\[
 C\min(\sigma_2-1,2-\sigma_2)>3+t|\sigma_1-\sigma_2|.
\]
Claim~\ref{c:binary-final-target} of Lemma~\ref{lem:binary-final-completion} gives a run-edit list $\Omega_*=(\omega_{*,1},\ldots,\omega_{*,n_*})$ from $X_*$
to a translated copy of $\operatorname{RunFactor}_{p_h+1,t+C}(V)$.
Its insertion count is $C$. Claim~\ref{c:binary-final-exchanges} of Lemma~\ref{lem:binary-final-completion} bounds its exchange count by $(t+C)^2$.

Let us assemble the run-edits on the full path.
Apply $\Omega_1,\ldots,\Omega_h$ and $\Omega_*$ successively to their factors in the whole current path.
Let $B^{\mathrm L}$ and $B^{\mathrm R}$ be the prefix and suffix paths of the current factor $Z$, so the current path is $B^{\mathrm L}\ccn Z\ccn B^{\mathrm R}$.
Replace a local insertion $\operatorname{Ins}(Z,i,R_j)$ by
\[
 \operatorname{Ins}(B^{\mathrm L}\ccn Z\ccn B^{\mathrm R},|B^{\mathrm L}|+i,R_j).
\]
Replace a local exchange $\operatorname{Exc}(Z;i,R_a,R_b)$ by
\[
 \operatorname{Exc}(B^{\mathrm L}\ccn Z\ccn B^{\mathrm R};|B^{\mathrm L}|+i,R_a,R_b).
\]
Induction on the local run-edits gives translated copies of $\operatorname{RunFactor}_{p_{s-1}+1,L+C+q_s}(V)$ for $s=1,\ldots,h$, followed by $\operatorname{RunFactor}_{p_h+1,t+C}(V)$.
Every intermediate path is binary and the insertion and exchange counts are unchanged.

Let us count the simple-run insertions and exchanges.
Let $Q$ be the resulting path and let $T_0=p_h+t+C$.
Then,
\[
\begin{gathered}
 Q=\operatorname{Trans}(\operatorname{RunFactor}_{1,T_0}(V),\operatorname{Start}_{\mathrm{pt}}(Q)),\\
 T_0=l+\sum_{s=1}^h(C+q_s)+C
 \le l+h\eta_0 L+C\le(1+\eta)l+B\le T.
\end{gathered}
\]
If $T_0<T$, append $\operatorname{RunFactor}_{T_0+1,T-T_0}(V)$.
Concatenating the run-edit lists on the whole path gives $\Omega$.
Let $\zeta=\operatorname{Start}_{\mathrm{pt}}(Q)$ and write $\Omega=(\omega_1,\ldots,\omega_n)$.
Define
\[
\begin{aligned}
 I_R&=\{s\in[n]_{\Z}:\omega_s\text{ is a simple-run insertion}\},\\
 I_X&=\{s\in[n]_{\Z}:\omega_s\text{ exchanges adjacent simple runs}\},\\
 K&=\#I_R,\qquad k=\#I_X.
\end{aligned}
\]
Then, $[n]_{\Z}=I_R\sqcup I_X$.
Every intermediate path is binary. Each exchange has the unequal-length property in Definition~\ref{def:binary-run-edit-lists}.
The endpoint construction and Definition~\ref{def:binary-run-edit-lists} give
\[
 \Omega\in\operatorname{BinEdit}_{K,k}
 \bigl(U,\operatorname{Trans}(\operatorname{RunFactor}_{1,T}(V),\zeta)\bigr).
\]

\prf{Claim~\ref{c:binary-mechanical-insertions}}
The simple-run insertion count is
\[
 K=(T_0-l)+(T-T_0)=T-l.
\]

\prf{Claim~\ref{c:binary-mechanical-exchanges}}
The exchange counts of $\Omega_1,\ldots,\Omega_h$ and $\Omega_*$ give
\[
 k\le7\eta_0 hL+(t+C)^2\le7\eta l+B.
\]
\end{proof}

\begin{proposition}\label{prop:exchange-composition}
Let $m=3$. Let $K\in\Z_{\ge0}$, $k\in\Z_{\ge0}$, and $P\in\Z_{\ge0}$ satisfy $P\ge K+4k$.
Let $U$ and $V$ be binary runs and fix $\Omega\in\operatorname{BinEdit}_{K,k}(U,V)$.
Let $u\in[2]_{\Z}$ and $v\in[2]_{\Z}$. Let $W_{\mathrm L}$ and $W_{\mathrm R}$ be binary runs with $P$ simple runs each.
Let $V_*=W_{\mathrm L}\ccn R_u\ccn V\ccn R_v\ccn W_{\mathrm R}$.
Then, the following statements hold.
\begin{enumerate}
\item\label{c:exchange-composition-map} There is a map
\[
 \Phi:\mathcal C_0(U)\longrightarrow\mathcal C_0(V_*)
\]
such that
\begin{equation}\label{eq:exchange-composition-fibers}
 \#\Phi^{-1}(\{\xi\})\le8^k\qquad(\xi\in\mathcal C_0(V_*)).
\end{equation}
\item\label{c:exchange-composition-count} We have
\begin{equation}\label{eq:exchange-composition}
 C_0(V_*)\ge\frac{C_0(U)}{8^k}.
\end{equation}
\end{enumerate}
\end{proposition}
\begin{proof}
\prf{Claim~\ref{c:exchange-composition-map}} Let us compose the maps for the individual run-edits.
Let $\Omega=(\omega_s)_{s\in[K+k]_{\Z}}$.
For $s\in[0,K+k]_{\Z}$, let $U_s=\operatorname{RunEResult}((\omega_1,\ldots,\omega_s),U)$.
Choose $I_R$, $I_X$, and $(i_s,j_s)_{s\in I_R}$ satisfying the conditions for $\Omega\in\operatorname{BinEdit}_{K,k}(U,V)$ in Definition~\ref{def:binary-run-edit-lists}.
Let us apply $\omega_1,\ldots,\omega_{K+k}$ to the copy of $U$ in $\operatorname{Shell}(U;u,v,0)$ and bound the surplus.
Define $a_j\in[2]_{\Z}$ and $b_j\in[2]_{\Z}$ for $j\in[P]_{\Z}$ by
\[
\begin{aligned}
 W_{\mathrm L}&=\operatorname{Trans}(R_{a_P}\ccn \cdots\ccn R_{a_1},\operatorname{Start}_{\mathrm{pt}}(W_{\mathrm L})),\\
 W_{\mathrm R}&=\operatorname{Trans}(R_{b_1}\ccn \cdots\ccn R_{b_P},\operatorname{Start}_{\mathrm{pt}}(W_{\mathrm R})).
\end{aligned}
\]
For $(f,\gamma)\in\mathcal C_0(U)$, let
\[
 Q_0=\operatorname{Shell}(U;u,v,0),\qquad
 (h_0,\eta_0)=\Phi_{\mathrm{sh}}(f,\gamma).
\]
Here $\Phi_{\mathrm{sh}}$ is the map in Claim~\ref{c:shell-recolor} of Lemma~\ref{lem:shell} with $j=1$ and $\epsilon=0$.
Claims~\ref{c:shell-connected} and~\ref{c:shell-residual} of Lemma~\ref{lem:shell} give $(h_0,\eta_0)\in\mathcal C_0(Q_0)$.
For $s\in I_R$, let $L\ccn U_{s-1}\ccn R$ be the current path and apply Lemma~\ref{lem:simple-run-insertion} at $|L|+i_s$.
Equation~\eqref{eq:run-insertion-surplus} shows that this run-edit adds one to the surplus.
For $s\in I_X$, let $e_s=1+\#(I_X\cap[s-1]_{\Z})$.
Let $Z=L\ccn U_{s-1}\ccn R$ be the current path and let $Z'=L\ccn U_s\ccn R$.
Choose
\[
 \Phi_s\in\operatorname{WrapExch}(Z,Z';a_{e_s},b_{e_s}).
\]
Lemma~\ref{lem:exchange} gives a member of this set because $Z'$ exchanges the two simple runs specified by $\omega_s$ inside $Z$.
The choice of $\Phi_s$ is independent of the input colored grading.
Apply $\Phi_s$ to the colored grading at step $s-1$.
Claim~\ref{c:exchange-residual} of Lemma~\ref{lem:exchange} bounds the increase in surplus by three.
Thus, the colored grading $(h,\eta)$ on the resulting path $Q$ satisfies
\[
 0\le\operatorname{Surp}(Q,h)\le K+3k\le P-k.
\]

Let us reduce the surplus to zero.
Apply Lemma~\ref{lem:surplus-reduction} with $n=P-k$ and
\[
 \alpha'=(a_{k+j})_{j\in[P-k]_{\Z}},\qquad
 \mu'=(b_{k+j})_{j\in[P-k]_{\Z}}.
\]
Lemma~\ref{lem:word-reindexing} translates the final colored grading to $V_*$.
This defines a map $\Phi\colon\mathcal C_0(U)\to\mathcal C_0(V_*)$.

Let us bound the fibers of $\Phi$.
The map $\Delta_{Q,\alpha',\mu'}$ from Lemma~\ref{lem:surplus-reduction} is injective.
Restriction recovers the colored grading before each simple-run insertion
by~\eqref{eq:run-insertion-grades} and~\eqref{eq:run-insertion-recolor}.

Let us count the choices at the $k$ exchanges.
Claim~\ref{c:exchange-fibers} of Lemma~\ref{lem:exchange} gives at most eight
inputs at each exchange.
Claim~\ref{c:shell-recolor} of Lemma~\ref{lem:shell} recovers $(f,\gamma)$ by restriction from $(h_0,\eta_0)$.
Recovering the inputs in the reverse order of $\Omega$ gives
\[
 \#\Phi^{-1}(\{\xi\})\le8^k\qquad(\xi\in\mathcal C_0(V_*)).
\]
This proves~\eqref{eq:exchange-composition-fibers}.

\prf{Claim~\ref{c:exchange-composition-count}} Let us sum the fiber bounds to compare the grading counts.
Summing~\eqref{eq:exchange-composition-fibers} over $\xi\in\mathcal C_0(V_*)$ gives
\[
 C_0(U)=\sum_{\xi\in\mathcal C_0(V_*)}\#\Phi^{-1}(\{\xi\})
 \le8^k C_0(V_*).
\]
Since $8^k>0$, division by $8^k$ proves~\eqref{eq:exchange-composition}.
\end{proof}

\begin{proposition}\label{prop:three-irrational}
We have the following.
\begin{enumerate}
\item\label{c:three-irrational-limit} For every irrational $r\in(1,2)$,
there is $\Upsilon_3(r)\in(0,\infty)$ such that
\[
 \lim_{\substack{a\to\infty,\ (a,b)\in\cA_3\\\frac{b}{a}\to r}}
 \frac{\log E_3(a,b)}a=\Upsilon_3(r).
\]
\item\label{c:three-irrational-continuity} Let $\Upsilon_3(1)=\log\left(\frac{256}{27}\right)$.
Use $\Upsilon_3(r)$ from~\eqref{eq:rational-period} for rational $r\in(1,\frac32]$ and from Claim~\ref{c:three-irrational-limit} for irrational $r\in(1,\frac32)$.
Then, $\Upsilon_3|_{[1,\frac32]}$ takes values in $(0,\infty)$ and satisfies
\[
 \lim_{\substack{s\to r\\s\in[1,\frac32]}}\Upsilon_3(s)=\Upsilon_3(r)
 \qquad\left(r\in[1,\frac32]\right).
\]
\end{enumerate}
\end{proposition}
\begin{proof}
\prf{Claim~\ref{c:three-irrational-limit}} Let us choose the comparison parameters for Lemma~\ref{lem:binary-mechanical} and Claim~\ref{c:exchange-composition-count} of Proposition~\ref{prop:exchange-composition}.
Consider $\eta$ with $0<\eta<\frac{1}{1024}$. Let $d=1-144\eta$.
Choose $\delta$ and $B$ from Lemma~\ref{lem:binary-mechanical} for $r$ and $\eta$.
Choose relatively prime $p\in\Z_{>0}$ and $q\in\Z_{>0}$ such that $q\ge2$,
$\alpha=\frac{p}{q}\in(1,2)$, and $|\alpha-r|<\delta$.
Let $V_*$ have $\nu$ simple runs and let $P=\floor{64\eta\nu}$.
Let us translate $V_*$ so that
\[
 V_*=R_{a_P}\ccn \cdots\ccn R_{a_1}\ccn R_u\ccn V\ccn R_v\ccn R_{b_1}\ccn \cdots\ccn R_{b_P},
 \qquad a_j,b_j,u,v\in[2]_{\Z}.
\]
Thus, $V$ has $T=\nu-2P-2$ simple runs.
For $l=d\nu+O(1)$, we have
\[
 T-(1+\eta)l=(15\eta+144\eta^2)\nu+O(1)\longrightarrow\infty.
\]
Hence, $T\ge(1+\eta)l+B$ for all sufficiently large $\nu$.

Let us specify the two applications of Lemma~\ref{lem:binary-mechanical}.
Consider $(a,b)\in\cA_3$ as $a\to\infty$ and $\frac ba\to r$.
\begin{enumerate}[label=(\roman*)]
\item For the lower bound, let
\[
\begin{gathered}
 \nu=a+1,\quad j=\floor{\frac{d\nu}{q}},\quad l=jq,\quad
 U=W_{p,q}^j,\quad V_*=\mathcal U(a,b),\\
 (\sigma_1,\sigma_2,u_1,u_2)
 =\bigl(\alpha,\Gamma(a,b),0,\{(P+1)\Gamma(a,b)\}\bigr).
\end{gathered}
\]
\item For the upper bound, let
\[
\begin{gathered}
 l=a+1,\quad N=\ceil{\frac{l}{dq}},\quad \nu=Nq,\quad
 U=\mathcal U(a,b),\quad V_*=W_{p,q}^N,\\
 (\sigma_1,\sigma_2,u_1,u_2)
 =\bigl(\Gamma(a,b),\alpha,0,\{(P+1)\alpha\}\bigr).
\end{gathered}
\]
\end{enumerate}

Let us verify admissibility for each of the two comparisons.
In each case, let $A_1,\ldots,A_l$ be the simple runs of $U$ in path order.
Equations~\eqref{eq:ambient}, \eqref{eq:frame-word}, and~\eqref{eq:rational-period} show that
$(l,T,\sigma_1,\sigma_2,u_1,u_2,(A_1,\ldots,A_l))$ is
$(2,r,\eta,\delta,B)$-admissible for all sufficiently large $a$.
The path $V$ is a translated copy of
$\operatorname{RunFactor}_{1,T}(\operatorname{MechRun}(\sigma_2,u_2))$.
In case (i), the final $R_1$ of $\mathcal U(a,b)$ belongs to $W_{\mathrm R}$.
In case (ii), we have $A_l=R_1$ and $|A_l|_N=1\in[2]_{\Z}$, as permitted by condition~\ref{hyp:mechanical-source} of Definition~\ref{def:admissible-mechanical}.

Let us bound the insertions and exchanges in each comparison.
Lemma~\ref{lem:binary-mechanical} gives an insertion count $K=T-l$ and an exchange count $k$ satisfying
\[
 \lim_{\nu\to\infty}\frac K\nu=16\eta,\qquad k\le7\eta l+B\le8\eta\nu,
 \qquad K+4k\le P
\]
for all sufficiently large $\nu$.
Apply Claim~\ref{c:exchange-composition-count} of Proposition~\ref{prop:exchange-composition} and Lemma~\ref{lem:word-reindexing} in each case.

First, let us obtain the lower bound from case (i), keeping $\eta$, $p$, and $q$ fixed.
Equation~\eqref{eq:dictionary} gives
\[
 \frac{\log E_3(a,b)}a
 \ge\frac ja\frac{\log C_0(W_{p,q}^j)}j
       -\frac ka\log8-\frac{\log3}a,
 \qquad
 \frac ja\longrightarrow\frac dq,\quad \frac\nu a\longrightarrow1.
\]

Second, let us obtain the upper bound from case (ii).
Equation~\eqref{eq:dictionary} gives
\[
 \frac{\log E_3(a,b)}a
 \le\frac Na\frac{\log C_0(W_{p,q}^N)}N
       +\frac ka\log8-\frac{\log3}a,
 \qquad
 \frac Na\longrightarrow\frac1{dq},\quad \frac\nu a\longrightarrow\frac1d.
\]

Let us apply the periodic limit to both bounds.
We have $j\to\infty$ in case (i) and $N\to\infty$ in case (ii).
Use $k\le8\eta\nu$ and Claim~\ref{c:periodic-limit} of Proposition~\ref{prop:periodic} to obtain
\begin{equation}\label{eq:three-irrational-comparison}
 \begin{aligned}
 \liminf_{\substack{a\to\infty,\ (a,b)\in\cA_3\\\frac{b}{a}\to r}}
 \frac{\log E_3(a,b)}a
 &\ge d\Upsilon_3(\alpha)-8\eta\log8,\\[1ex]
 \limsup_{\substack{a\to\infty,\ (a,b)\in\cA_3\\\frac{b}{a}\to r}}
 \frac{\log E_3(a,b)}a
  &\le\frac{\Upsilon_3(\alpha)+8\eta\log8}{d}.
 \end{aligned}
\end{equation}
For each fixed $\eta$, we choose $\alpha=\frac{p}{q}$, $\delta$, and $B$
before letting $a\to\infty$ and $\frac ba\to r$ through $(a,b)\in\cA_3$.

Let us bound the difference between the upper and lower bounds in~\eqref{eq:three-irrational-comparison}. Lemma~\ref{lem:word-bounds} with $m=3$ and $L=2$, together with Claim~\ref{c:periodic-limit} of Proposition~\ref{prop:periodic} and~\eqref{eq:rational-period}, gives
\[
 \Upsilon_3(\alpha)\le3\log8.
\]
For each fixed $\eta$, equation~\eqref{eq:three-irrational-comparison} bounds the lower and upper limits for every sequence $((a_n,b_n))_{n\ge1}$ in $\cA_3$ satisfying
\[
 \lim_{n\to\infty}a_n=\infty,\qquad
 \lim_{n\to\infty}\frac{b_n}{a_n}=r.
\]

Let us prove that the upper bound minus the lower bound in~\eqref{eq:three-irrational-comparison} tends to zero as $\eta\downarrow0$.
The difference is at most
\[
 3\log8\left(\frac1d-d\right)
 +8\eta\log8\left(1+\frac1d\right)
 \longrightarrow0.
\]
Since $r\in(1,2)\subset I_3$, Lemma~\ref{lem:approaching-sequence} gives a sequence $((a_n,b_n))_{n\ge1}$ in $\cA_3$ satisfying
\[
 \lim_{n\to\infty}a_n=\infty,\qquad
 \lim_{n\to\infty}\frac{b_n}{a_n}=r.
\]
For any two such sequences, denote their lower and upper limits by $L_i$ and $U_i$, $i\in[2]_{\Z}$.
They are finite by Lemma~\ref{lem:word-bounds} and~\eqref{eq:dictionary}.
The common interval in~\eqref{eq:three-irrational-comparison} gives
\[
 U_i-L_j\le3\log8\left(\frac1d-d\right)
             +8\eta\log8\left(1+\frac1d\right)
 \qquad(i,j\in[2]_{\Z}).
\]
Letting $\eta\downarrow0$ gives $U_i\le L_j$ for all $i,j$.
Since $L_i\le U_i$, both sequences converge and their limits are equal.
Define $\Upsilon_3(r)$ to be this common limit.

Let us check positivity of the limit.
The run $\mathcal U(a,b)$ has $a+1$ simple runs of vertical length at most two for all sufficiently large $a$.
Lemma~\ref{lem:word-bounds} and~\eqref{eq:dictionary} give
\[
 \frac{\log E_3(a,b)}a
 \ge\frac{\lceil\frac{a+1}{2}\rceil-1}{a}\log3
 \longrightarrow\frac12\log3.
\]
Hence, $\Upsilon_3(r)\ge\frac12\log3>0$.

\prf{Claim~\ref{c:three-irrational-continuity}} Let us combine the three kinds of slopes before applying the continuity criterion.
Claim~\ref{c:three-irrational-limit}, Theorem~\ref{thm:rational}, and Corollary~\ref{cor:central} give, respectively, the irrational, noninteger rational $r>1$, and $r=1$ cases of
\[
 \lim_{\substack{a\to\infty,\ (a,b)\in\cA_3\\\frac ba\to r}}
 \frac{\log E_3(a,b)}a=\Upsilon_3(r)
 \qquad\left(r\in[1,\frac32]\right).
\]
For rational $r\in(1,\frac32]$, Claim~\ref{c:periodic-limit} of Proposition~\ref{prop:periodic} and~\eqref{eq:rational-period} give $\Upsilon_3(r)>0$.
At $r=1$, we have $\Upsilon_3(1)=\log\frac{256}{27}>0$.
Lemma~\ref{lem:approaching-sequence} gives hypothesis~\ref{hyp:approaching-sequences} of Lemma~\ref{lem:continuity} for $J=[1,\frac32]$.
Claim~\ref{c:continuity-function} of Lemma~\ref{lem:continuity} therefore gives continuity on $[1,\frac32]$.
Claim~\ref{c:continuity-endpoint-function} of Lemma~\ref{lem:continuity} specifies the endpoint limits:
\[
 \lim_{s\downarrow1}\Upsilon_3(s)=\Upsilon_3(1),\qquad
 \lim_{s\uparrow\frac32}\Upsilon_3(s)=\Upsilon_3\left(\frac32\right)
 \qquad\left(s\in\left[1,\frac32\right]\right).
\]
\end{proof}


